\documentclass[12pt]{article}


\usepackage{macros}
\usepackage{macros-gen-ro}
\usepackage{macros-math}
\usepackage{macros-logic}

\newcommand{\seminarno}{0}
\newcommand{\curs}{Logic\u a matematic\u a}
\newcommand{\departamentan}{ FMI, Mate, Anul I}


\newsavebox{\descrierecurs}
\savebox{\descrierecurs}{\raisebox{-11pt}{\parbox[t]{8cm} {\bf \departamentan \\ \curs}}}

\newcommand{\titlu}{ 
\usebox{\descrierecurs}  
%\bigskip
%\begin{flushright} \datum \end{flushright}
\vspace{10mm}
\begin{center}{\Large \bf 
Seminar \seminarno}
 \end{center} \vspace{5mm}
}


\newcounter{exercitiuno}
\newcommand{\semex}{\stepcounter{exercitiuno}{\vspace{0.3cm}\noindent \bf (S\seminarno.\arabic{exercitiuno}) }}

\begin{document}


\renewcommand{\seminarno}{1}
\setcounter{exercitiuno}{0}

\titlu


\renewcommand{\solution}[1]{}
%\renewcommand{\details}{}


\parindent0pt



\semex
Fie $A,B,C$ mul\c{t}imi. Demonstra\c{t}i c\u{a}
\be
\item $A\times (B\cup C) = (A\times B)\cup (A\times C)$.
\item $A\times (B\cap C) = (A\times B)\cap (A\times C)$. 
\ee
\solution{
\be
\item 
Avem c\u a  \\[1mm]
\bt{lll} 
$x \in A\times (B\cup C)$  & \ddacam &  exist\u a $ a \in A$ \c si $y \in B \cup C$ cu $x=(a,y)$ \\
&\ddacam &  exist\u a  $ a \in A$ \c si $y \in B$ sau $y \in C$ cu $x=(a,y)$ \\
&\ddacam & \big(exist\u a  $a \in A$ \c si $y \in B$ cu $x=(a,y)$\big) \\
&& sau \big(exist\u a  $a \in A$ \c si $y \in C$ cu $x=(a,y)$\big) \\
&\ddacam & $x \in A \times B$ sau $x \in A \times C$ \\
&\ddacam & $x \in (A\times B)\cup (A\times C)$.
\et

Prin urmare, $A\times (B\cup C) =  (A\times B)\cup (A\times C)$. 

\item Demonstr\u am prin dubl\u a  incluziune. 

$\se$ Fie $x \in A\times (B\cap C)$. Atunci exist\u a  $a \in A$ \c si $y \in B \cap C$ cu $x=(a,y)$. Deoarece $y\in B$, rezult\u a c\u a $x\in A\times B$. Deoarece $y\in C$, rezult\u a c\u a $x\in A\times C$.  Prin urmare, $x\in  (A\times B)\cap (A\times C)$.

$\supseteq$ Fie $x \in (A\times B)\cap (A\times C)$. Atunci $x\in A\times B$ \c si $x\in A\times C$, a\c sadar, $x=(a_1,y_1)$ cu $a_1\in A, y_1\in B$ \c si
$x=(a_2,y_2)$ cu $a_2\in A, y_2\in B$. Deoarece $x=(a_1,y_1)=(a_2,y_2)$, trebuie s\u a avem $a_1=a_2$ \c si $y_1=y_2$. Fie $a:=y_1=y_2$ \c si $y:=y_1=y_2$.
Atunci $x=(a,y)$, cu $a\in A$ \c si $y\in B\cap C$. Rezult\u a c\u a $x\in A\times (B\cap C)$.

Am demonstrat c\u a  $A\times (B\cap C) = (A\times B)\cap (A\times C)$.
\ee
}


\semex 
Fie $A=\{a,b,c,d\}$ \c{s}i $R=\{(a,b),(a,c),(c,d),(a,a),(b,a)\}$ o rela\c{t}ie binar\u{a} pe $A$. Care este compunerea $R\circ R$? Care este inversa $R^{-1}$ a lui $R$?  \\
\solution{Ob\c{t}inem 
\bua R \circ R &=& \{(a,a), (a,b), (a,c), (a,d), (b,a), (b,b), (b,c)\},\\
R^{-1} &=& \{ (a,a), (a,b), (b,a), (c,a), (d,c)\}.
\eua
}

\semex
S\u a se demonstreze asociativitatea compunerii rela\c{t}iilor. \\
\solution{
Fie $R \subseteq A \times B$, $Q \subseteq B \times C$, $S \subseteq C \times D$ trei rela\c tii. Vrem  s\u a demonstr\u am c\u a 
$$(R \circ Q) \circ S = R \circ (Q \circ S).$$ 

$\se$ Fie $(a,d) \in (R \circ Q) \circ S$. Atunci exist\u a $c \in C$ cu $(a,c) \in R \circ Q$ \c si $(c,d) \in S$. Din faptul c\u a $(a,c) \in R \circ Q$ avem c\u a exist\u a $b \in B$ cu $(a,b)\in R$ \c si $(b,c) \in Q$. Din faptul c\u a $(b,c) \in Q$ \c si $(c,d) \in S$, avem $(b,d) \in Q \circ S$. 
Am ob\c tinut c\u a $(a,b)\in R$ \c si $(b,d) \in Q \circ S$. Prin urmare, $(a,d) \in R \circ (Q \circ S)$. 

$\supseteq$ Se demonstreaz\u a analog. 
}

\semex
Fie $(A,\leq)$ o mul\c{t}ime par\c{t}ial ordonat\u{a} \c{s}i  $\emptyset \ne S\se A$. 
Demonstra\c ti urm\u atoarele:
\be
\item At\^{a}t minimul, c\^{a}t \c{s}i maximul lui $S$ sunt unice (dac\u{a} exist\u{a}).
\item Dac\u a $\min S$ exist\u a, atunci $\min S$ este element minimal al lui $S$.
\item Dac\u a $\max S$ exist\u a, atunci $\max S$ este element maximal al lui $S$. 
\ee
\solution{
\be
\item Presupunem c\u a $x$ \c si $x'$ sunt minime ale lui $S$ \c si demonstr\u am c\u a sunt egale.

Deoarece  $x$ este minim al lui $S$, avem c\u a  $x \leq y$  pentru orice $y \in S$. Lu\^and $y:=x'$ ob\c tinem
\beq
x \leq x' \label{xx'}
\eeq

Deoarece  $x'$ este minim al lui $S$, avem c\u a  $x' \leq y$ pentru orice $y \in S$. Lu\^and $y:=x$ ob\c tinem
\beq
x' \leq x \label{x'x}
\eeq
Din \eqref{xx'} \c si \eqref{x'x}, folosind faptul c\u a $\leq $ este antisimetric\u a, rezult\u a c\u a $x=x'$.

Se procedeaz\u a asem\u an\u ator pentru maxim.
\item Fie $x = \min S$.  Fie $t\in S$ \ai  $t\leq x$. Deoarece $x$ este minimul lui $S$, avem \c si  c\u a $x\leq t$. Din $t\leq x$, $x\leq t$ \c si antisimetria lui
$\leq$ rezult\u a c\u a $x=t$. Prin urmare, $x$ este element minimal al lui $S$.

\item Se demonstreaz\u a similar. 
\ee
}



\semex
Fie $X$ o mul\c time. S\u a se arate c\u a nu exist\u a o func\c tie surjectiv\u a cu domeniul $X$ \c si codomeniul $\cP(X)$. \\
\solution
{
Presupunem c\u a ar exista \c si fie $f: X \to \cP(X)$ surjectiv\u a. Fie mul\c timea 
$$A = \{x \in X \mid x \notin f(x) \} \in \cP(X).$$
 Dat fiind c\u a $f$ este surjectiv\u a, exist\u a $y \in X$ cu $f(y)=A$. Dar atunci:
$y \in A \Leftrightarrow y \notin f(y) = A \Leftrightarrow y \notin A$
ceea ce este o contradic\c tie.
}

\semex
Fie $\mathcal{F} = \{B \subseteq \N \mid B \text{ este finită} \}$.  Să se demonstreze că $(\mathcal{F}, \subseteq)$ nu este inductiv ordonată. \\
\solution{
Consider\u am submul\c timea 
\bce 
$\cC = \{ A_n \mid n\in \N \}$, unde $A_n = \{ 0,1, 2, \ldots, n \}$.
\ece 
Se arat\u a u\c sor c\u a  $\cC$ este total ordonată: pentru orice $m, n\in \N$, avem c\u a $A_m \se A_n$ dacă $m \leq n$  \c si $A_n \se A_m$ dacă $n \leq m$. 

Presupunem c\u a $B\in \mathcal{F}$ este majorant al lui $\cC$. Deoarece $B$ este finit\u a, exist\u a $M=\max B$. Rezult\u a c\u a $B \se A_M  \subset A_{M+1}$.  Am ob\c tinut o 
contradic\c tie.  Prin urmare, $\cC$ nu are majoran\c ti. 

Deci, $(\mathcal{F}, \subseteq)$ nu este inductiv ordonată.
}

\end{document}



